\setauthor{\firstauthor}

% ============================================================
%  Authentifizierung & Autorisierung -- Implementierung
%  (gehört in Kapitel 3: Praktische Umsetzung)
% ============================================================

Dieses Kapitel beschreibt die konkrete Implementierung der
Authentifizierung und Autorisierung in LearnHub -- aufbauend auf den
Grundlagen aus Kapitel~\ref{chapter:tech}.

% ----------------------------------------------------------
\subsection{Frontend}
% ----------------------------------------------------------

\subsubsection{Keycloak-Initialisierung}

Das Frontend nutzt die Bibliothek \texttt{keycloak-angular} für die
Integration.~\cite{keycloak-docs} Die Initialisierung erfolgt beim
App-Start über einen \texttt{APP\_INITIALIZER}, der sicherstellt,
dass Keycloak vollständig geladen ist, bevor die Anwendung
startet:

\begin{lstlisting}[style=tsstyle,
                   caption={Keycloak-Initialisierung (keycloak.provider.ts)},
                   label={lst:keycloak-init}]
export function initializeKeycloak(
  keycloak: KeycloakService,
  injector: Injector
) {
  return async () => {
    const authenticated = await keycloak.init({
      config: {
        url: 'https://auth.htl-leonding.ac.at',
        realm: '2526_5bhitm',
        clientId: 'frontend',
      },
      initOptions: {
        onLoad: 'login-required',
        checkLoginIframe: false,
        pkceMethod: 'S256',
      },
      enableBearerInterceptor: false,
    });

    if (authenticated) {
      const userInitService =
        injector.get(UserInitializationService);
      await userInitService.initializeUser();
    }

    return authenticated;
  };
}
\end{lstlisting}

Die Option \texttt{onLoad: 'login-required'} bewirkt, dass nicht
angemeldete Nutzer:innen sofort zur Keycloak-Loginseite
weitergeleitet werden. \texttt{pkceMethod: 'S256'} aktiviert den
Authorization Code Flow mit PKCE.~\cite{oauth2-rfc} Nach erfolgreicher
Anmeldung ruft \texttt{UserInitializationService.initializeUser()}
den Backend-Endpunkt auf, der den Nutzer anlegt oder abruft.

\subsubsection{Auth Guard}

Der \texttt{AuthGuard} schützt Angular-Routen vor unauthentifizierten
Zugriffen und prüft optional auf Rollenanforderungen:

\begin{lstlisting}[style=tsstyle,
                   caption={AuthGuard (auth.guard.ts)},
                   label={lst:auth-guard}]
@Injectable({ providedIn: 'root' })
export class AuthGuard extends KeycloakAuthGuard {

  public async isAccessAllowed(
    route: ActivatedRouteSnapshot,
    state: RouterStateSnapshot
  ): Promise<boolean> {

    if (!this.authenticated) {
      await this.keycloak.login({
        redirectUri: window.location.origin + state.url,
      });
      return false;
    }

    const requiredRoles = route.data['roles'] as string[];
    if (!Array.isArray(requiredRoles) || requiredRoles.length === 0) {
      return true;
    }

    const hasRole = requiredRoles.some(r => this.roles.includes(r));
    if (!hasRole) {
      this.router.navigate(['/not-authorized']);
      return false;
    }

    return true;
  }
}
\end{lstlisting}

Ist der Nutzer nicht angemeldet, erfolgt eine automatische
Weiterleitung zur Keycloak-Loginseite. Fehlt eine erforderliche Rolle,
wird zur \texttt{/not-authorized}-Seite weitergeleitet. Routen ohne
\texttt{roles}-Angabe sind für alle angemeldeten Nutzer:innen
zugänglich.

\subsubsection{Auth Interceptor}

Der \texttt{AuthInterceptor} fügt jeder ausgehenden HTTP-Anfrage
automatisch das Bearer-Token als Header hinzu:

\begin{lstlisting}[style=tsstyle,
                   caption={Auth Interceptor (auth.interceptor.ts)},
                   label={lst:auth-interceptor}]
export const authInterceptor: HttpInterceptorFn = (req, next) => {
  const keycloak = inject(KeycloakService);
  const token = keycloak.getKeycloakInstance().token;

  if (token) {
    const cloned = req.clone({
      setHeaders: { Authorization: `Bearer ${token}` }
    });
    return next(cloned).pipe(
      catchError((error: HttpErrorResponse) => {
        if (error.status === 401) {
          keycloak.login();
        }
        return throwError(() => error);
      })
    );
  }

  keycloak.login();
  return throwError(() => new Error('No authentication token'));
};
\end{lstlisting}

Bei einem HTTP-401-Fehler (abgelaufenes Token) wird der Nutzer
automatisch erneut zur Loginseite weitergeleitet.

\subsubsection{User Initialization Service}

Nach erfolgreicher Anmeldung registriert \texttt{UserInitializationService}
den Nutzer im Backend und hält das Ergebnis als reaktiven
\texttt{BehaviorSubject} für alle Komponenten verfügbar:

\begin{lstlisting}[style=tsstyle,
                   caption={UserInitializationService (user-initialization.service.ts)},
                   label={lst:user-init-service}]
@Injectable({ providedIn: 'root' })
export class UserInitializationService {
  private currentUserSubject =
    new BehaviorSubject<UserSlim | null>(null);
  public currentUser$ = this.currentUserSubject.asObservable();

  async initializeUser(): Promise<UserSlim | null> {
    if (this.currentUserSubject.value) {
      return this.currentUserSubject.value;
    }
    if (!this.keycloakService.isLoggedIn()) return null;

    const user =
      await this.userService.registerOrGetCurrentUserAsync();
    this.currentUserSubject.next(user);
    return user;
  }
}
\end{lstlisting}

% ----------------------------------------------------------
\subsection{Backend}
% ----------------------------------------------------------

\subsubsection{JWT-Validierung}

Das Backend validiert bei jeder Anfrage das mitgeschickte JWT über
einen \texttt{JwtRequestFilter}. Der Filter prüft:

\begin{enumerate}
    \item Ist ein \texttt{Authorization: Bearer}-Header vorhanden?
    \item Ist die RS256-Signatur mit dem öffentlichen Keycloak-Schlüssel
    gültig?
    \item Ist das Token noch nicht abgelaufen?
\end{enumerate}

Bei ungültigem Token wird die Anfrage mit HTTP~401 abgebrochen. Ist
das Token gültig, wird ein \texttt{CustomSecurityContext} erstellt,
der die extrahierten Nutzerinformationen für alle Endpoints
bereitstellt:

\begin{lstlisting}[language=Java,
                   caption={CustomSecurityContext (CustomSecurityContext.java)},
                   label={lst:custom-security-context}]
public record CustomSecurityContext(
    String keycloakSub,
    String username,
    String email,
    String distinguishedName,
    boolean isStudent
) implements SecurityContext {

    @Override
    public Principal getUserPrincipal() {
        return () -> username;
    }

    @Override
    public boolean isUserInRole(String role) {
        return switch (role) {
            case "teacher" -> !isStudent;
            case "student" -> isStudent;
            default        -> false;
        };
    }

    @Override public boolean isSecure() { return true; }

    @Override
    public String getAuthenticationScheme() { return "Bearer"; }
}
\end{lstlisting}

\subsubsection{Rechte im Backend}

Da LearnHub einen eigenen \texttt{JwtRequestFilter} implementiert (kein
Standard-Quarkus-OIDC), verhält sich \texttt{@PermitAll} anders als im
Standard-Fall. Der Filter liest bei \texttt{@PermitAll}-Endpoints das
Token, falls vorhanden, \textit{optional}: Ist ein gültiges Token im
Header, wird der \texttt{CustomSecurityContext} befüllt. Ist keines
vorhanden, läuft die Anfrage trotzdem durch.~\cite{quarkus-security-jwt}

In der Praxis bedeutet das für LearnHub:

\begin{itemize}
    \item \texttt{@PermitAll} -- Token ist \textit{optional}. Ist eines
    vorhanden und gültig, steht der SecurityContext bereit. Wird in
    LearnHub für den \texttt{/register}-Endpoint verwendet, der das
    Token zwar verarbeitet, aber nicht erzwingt -- damit der
    Registrierungsaufruf direkt nach dem Login funktioniert, bevor
    der erste vollständige Anfragezyklus abgeschlossen ist.
    \item \texttt{@RolesAllowed("teacher")} -- Token ist
    \textit{Pflicht}, und der Nutzer muss als Lehrperson erkannt
    werden (kein \texttt{OU=Students} im DN). Schüler:innen erhalten
    HTTP~403.
    \item \texttt{@RolesAllowed(\{"student", "teacher"\})} -- Token
    ist Pflicht. Alle angemeldeten Nutzer:innen dürfen zugreifen,
    unabhängig von der Rolle.
\end{itemize}

\begin{lstlisting}[language=Java,
                   caption={Annotationen an Beispiel-Endpoints (QuestionResource.java)},
                   label={lst:roles-backend}]
@Path("/api/questions")
public class QuestionResource {

    // Alle angemeldeten Nutzer koennen oeffentliche Fragen abrufen
    @GET
    @RolesAllowed({"student", "teacher"})
    public List<QuestionDto> getPublicQuestions() { ... }

    // Nur Lehrpersonen duerfen Fragen freigeben
    @PATCH
    @Path("/{id}/public")
    @RolesAllowed("teacher")
    public Response approveQuestion(@PathParam("id") Long id,
        @Context SecurityContext sc) { ... }

    // Alle angemeldeten Nutzer koennen eigene Fragen anlegen
    @POST
    @RolesAllowed({"student", "teacher"})
    public Response createQuestion(QuestionDto dto,
        @Context SecurityContext sc) { ... }
}
\end{lstlisting}

\subsubsection{Testing: HTTP-Dateien und Swagger UI}

Da alle Endpoints ein gültiges JWT erfordern, wurde das Testen ohne
laufendes Frontend über zwei Werkzeuge gelöst.

\textbf{HTTP-Dateien} sind Textdateien mit der Endung \texttt{.http},
die direkt in IntelliJ IDEA ausgeführt werden können. Listing
\ref{lst:auth-http} zeigt die Datei für die Authentifizierung: Das
Token wird von Keycloak geholt und automatisch als Variable
gespeichert, die alle folgenden Requests verwenden.

\begin{lstlisting}[caption={auth.http -- Token holen und Endpoints testen},
                   label={lst:auth-http}]
### Token von Keycloak holen
POST https://auth.htl-leonding.ac.at/realms/2526_5bhitm
     /protocol/openid-connect/token
Content-Type: application/x-www-form-urlencoded

grant_type=password&client_id=frontend
  &username={{KC_USER}}&password={{KC_PASSWORD}}

> {%
    client.global.set("access_token",
      response.body.access_token);
  %}

### User registrieren / abrufen
POST http://localhost:8080/api/users/register
Authorization: Bearer {{access_token}}
\end{lstlisting}

Die Zugangsdaten werden in \texttt{http-client.private.env.json}
gespeichert, die zwingend in \texttt{.gitignore} eingetragen sein
muss:

\begin{lstlisting}[language=json,
                   caption={http-client.private.env.json (nicht ins Repository)},
                   label={lst:http-client-env}]
{
  "dev": {
    "KC_USER":     "your-htl-username",
    "KC_PASSWORD": "your-htl-password"
  }
}
\end{lstlisting}

\textbf{Swagger UI:} Quarkus generiert unter \texttt{/q/swagger-ui}
eine interaktive API-Oberfläche.~\cite{quarkus-security-jwt} Das
Backend definiert dafür ein \texttt{BearerAuth}-Sicherheitsschema in
der \texttt{application.properties}. Um Endpoints in Swagger zu
testen, wird zuerst über die HTTP-Datei ein Token von Keycloak geholt.
Anschließend öffnet man Swagger, klickt oben rechts auf
\textit{Authorize} und trägt den Token ein -- \textit{ohne} das Prefix
\texttt{Bearer}, da Swagger dieses automatisch hinzufügt. Nach dem
Bestätigen sind alle Endpoints mit dem Token authentifiziert und
können direkt im Browser ausgeführt werden.

Eine Besonderheit beim Testen war, dass das Team aufgrund des fehlenden
Zugriffs auf die Keycloak-Konfiguration keine eigenen Testkonten anlegen
konnte. Das Testing erforderte daher stets gültige HTL-Accounts, was
vor allem das automatisierte Testen von Authentifizierungsabläufen
erschwerte. Als Workaround wurden die HTTP-Dateien mit den eigenen
HTL-Zugangsdaten verwendet, was für manuelle Tests ausreichend war.

\subsubsection{User-Registrierung beim ersten Login}

Da Keycloak lediglich die Authentifizierung übernimmt, führt das
Backend eine eigene Nutzer-Tabelle. Diese ist notwendig, weil
Datenbankbeziehungen eine lokale ID erfordern und das Attribut
\texttt{isTeacher} persistent gespeichert werden muss. Beim ersten
Login wird der Nutzer automatisch angelegt:

\begin{lstlisting}[language=Java,
                   caption={User-Registrierung beim ersten Login (UserResource.java)},
                   label={lst:user-register}]
@POST
@Path("/register")
@Transactional
public Response registerOrGetUser(
    @Context SecurityContext securityContext
) {
    CustomSecurityContext ctx =
        (CustomSecurityContext) securityContext;

    return userService
        .findByKeycloakSub(ctx.keycloakSub())
        .map(existing -> Response.ok(existing).build())
        .orElseGet(() -> {
            User user    = new User();
            user.keycloakSub = ctx.keycloakSub();
            user.username    = ctx.username();
            user.email       = ctx.email();
            user.isTeacher   = !ctx.isStudent();
            userService.persist(user);
            return Response.status(Response.Status.CREATED)
                           .entity(user).build();
        });
}
\end{lstlisting}

Die Rollenzuweisung erfolgt automatisch anhand des
\texttt{distinguishedName} aus dem JWT: Enthält der DN den Eintrag
\texttt{OU=Students}, wird \texttt{isTeacher = false} gesetzt;
andernfalls \texttt{isTeacher = true}. Eine manuelle Zuweisung
ist nicht erforderlich.
